\section{Discussion}
\label{sec:discussion}

The exploration of plant microbiomes within extraterrestrial habitats represents a vital frontier for space agricultural research \citep{fu2016how}. Traditional spaceflight plant research has predominantly viewed seedlings through a reductionist lens, attributing physiological deviations solely to direct host cellular responses to microgravity and cosmic radiation \citep{paul2017genetic,choi2019variation}. By mining the non-host metatranscriptome from archival NASA OSDR accessions, this study provides compelling proof-of-concept that \textit{Arabidopsis} seedlings in space host complex, dynamic microbial consortia that cannot be ignored when evaluating spaceflight plant adaptation.

\subsection{The Seed as a Vehicle for Microbial Transfer to Space}
Standard spaceflight protocols mandate rigorous surface-sterilization of seeds using hypochlorite or ethanol solutions prior to plating \citep{dixit2021persistence}. The widespread survival and metabolic transcription of taxa like \textit{Ralstonia pickettii}, \textit{Pseudomonas fluorescens}, \textit{Herbaspirillum huttiense}, and \textit{Weizmannia ginsengihumi} confirm that surface chemical treatments fail to eliminate endophytic bacteria residing within the seed coat, aleurone layer, or embryo tissues.

This has profound implications for planetary protection and space farm biosecurity. As demonstrated by the persistent dominance of \textit{R. pickettii} across OSD-37 (up to 100\% in Cvi-0), single oligotrophic generalists can monopolize the seedling microbiome within closed containment hardware like the BRIC canisters. \textit{R. pickettii} is capable of surviving in extreme nutrient-depleted conditions, forming robust biofilms, and tolerating toxic metals and disinfectants \citep{segata2012metagenomic}. In microgravity, where fluid convection is absent and boundary-layer gas exchange is severely hindered, biofilm-forming bacteria could either exacerbate anoxic root-zone stress or, conversely, create localized microenvironments that buffer the plant against oxidative shock \citep{barker2023meta}.

\subsection{Discovery of \textit{Weizmannia ginsengihumi}: A Candidate Spaceflight Probiotic}
The detection of \textit{Weizmannia ginsengihumi} in OSD-38 highlights the power of archival metatranscriptome mining for bio-prospecting. In terrestrial agro-ecosystems, \textit{W. ginsengihumi} enhances host antioxidant defenses, synthesizes ginsenoside-modifying enzymes, and stimulates root elongation under abiotic stress \citep{kim2020weizmannia}. Given that plants in spaceflight chronically upregulate reactive oxygen species (ROS) scavengers, peroxidases, and oxidative stress pathways \citep{choi2019variation,kruse2020spaceflight}, the presence of an endophyte with intrinsic antioxidant-producing capabilities could represent a natural symbiotic adaptation. Future experiments should test whether intentional inoculation of crop seeds with \textit{W. ginsengihumi} or \textit{P. fluorescens} confers enhanced vigor, root architecture, and pathogen defense in orbital and lunar growth chambers.

\subsection{Human-Associated Commensals and Habitat Biosecurity}
The recurrent recovery of human skin commensals (\textit{Cutibacterium acnes}, \textit{Malassezia restricta}) and opportunistic bacteria (\textit{Escherichia coli}, \textit{Moraxella osloensis}) across OSD-37, OSD-120, and OSD-217 exposes the permeability of spaceflight hardware to crew microbiomes. Recent investigations by \citet{totsline2023microgravity,totsline2024simulated} demonstrated that microgravity induces wide stomatal opening while suppressing plant pattern-triggered immunity (PTI), facilitating stomatal ingression and internalization of human pathogens like \textit{Salmonella enterica} and \textit{E. coli} in salad crops. Our detection of active \textit{E. coli} and skin commensal transcripts in orbital plant RNA emphasizes the urgent need for stringent touch-transfer controls and microbiome monitoring in future crewed food production systems.

\subsection{Methodological Strengths and Limitations}
Extracting microbial signals from host RNA-seq data offers an exceptional, zero-mass, retroactive approach to space biology. Thousands of archival samples across OSDR and international repositories can be re-interrogated without necessitating new spaceflight missions. Nonetheless, several caveats must be considered:
\begin{enumerate}
    \item \textbf{Sequencing Library Bias:} Most OSDR plant RNA-seq libraries utilize poly(A) enrichment, which selectively enriches for eukaryotic mRNA. Non-polyadenylated bacterial transcripts are captured only through non-specific priming, partial polyadenylation during degradation, or co-purification. Consequently, relative read counts reflect a mixture of absolute abundance and transcriptional turnover.
    \item \textbf{Low Biomass Artifacts:} When host tissue biomass is minuscule—as in the micro-dissected root tips of OSD-120—the ratio of microbial reads to background noise drops below detection thresholds, as observed here.
    \item \textbf{Cross-Study Heterogeneity:} Differences in RNA preservation reagents (RNAlater vs.\ liquid nitrogen) significantly alter microbial RNA degradation kinetics \citep{kruse2017transcriptome}.
\end{enumerate}

To overcome these limitations, future spaceflight multi-omics initiatives should systematically incorporate dual RNA-seq protocols (total RNA extraction with concurrent ribosomal RNA depletion of both plant and microbial rRNAs), coupled with long-read metagenomic sequencing on orbit via MinION or Oxford Nanopore sequencers.
